% !TeX program = pdflatex
\documentclass[11pt,a4paper]{article}
\usepackage{../../recursos/latex/estilo-notas}

\begin{document}
\cabecalhoaula{05}{Métodos Não Paramétricos: Aspectos Teóricos}{AME §4.13 e cap.~5}

%==============================================================================
\section*{Objetivos da aula}
%==============================================================================
Ao final desta aula o aluno deve ser capaz de:
\begin{itemize}[leftmargin=1.4cm]
  \item entender o que é uma \emph{taxa de convergência} e por que ela mede a
        qualidade de um estimador não paramétrico;
  \item enunciar as hipóteses de \emph{suavidade} (Lipschitz, Sobolev) e seu papel;
  \item deduzir a taxa de convergência do \textbf{KNN} via decomposição
        viés--variância;
  \item enunciar a taxa de \emph{séries ortogonais} e a noção de \emph{taxa
        minimax};
  \item explicar a \textbf{maldição da dimensionalidade} e sua consequência sobre o
        tamanho de amostra necessário;
  \item descrever as hipóteses de \emph{esparsidade} e \emph{redundância} e quais
        métodos se adaptam a elas.
\end{itemize}

\begin{emsala}
Esta aula responde \emph{por que} os métodos da Aula~04 funcionam --- e quando
\emph{não} funcionam. Frase-guia (Kant): \emph{``experiência sem teoria é cega.''}
Pergunta de abertura, retomando o exemplo da Amazon (Aula~04): \emph{``por que o
KNN foi tão pior que o Lasso/XGBoost naquele problema de texto?''} A resposta é
toda teórica.
\end{emsala}

%==============================================================================
\section{Por que estudar taxas de convergência?}
%==============================================================================
Queremos comparar estimadores não pela performance em \emph{um} conjunto de dados,
mas pela rapidez com que o risco $\E[(\rhat(\X)-r(\X))^2]$ tende a zero conforme
$n\to\infty$. Uma \textbf{taxa de convergência} expressa esse decaimento em função
de $n$ (e da dimensão $p$, da suavidade etc.). Por exemplo, uma taxa $n^{-1/2}$ é
mais rápida (melhor) que $n^{-1/5}$.

\begin{ideia}
Estimação não paramétrica só é possível sob \textbf{suposições de suavidade}: se
$r$ pudesse oscilar arbitrariamente, nenhuma quantidade finita de dados a
determinaria. Suavidade é o que permite ``emprestar força'' de observações
vizinhas.
\end{ideia}

\noindent Duas medidas de suavidade que usaremos:
\begin{description}[leftmargin=2.2cm,style=nextline]
  \item[Lipschitz] $r$ é $L$-Lipschitz se
        $\abs{r(\x)-r(\x')}\le L\norm{\x-\x'}$ para todos $\x,\x'$. Controla
        \emph{quão rápido} $r$ pode variar.
  \item[Sobolev] $r$ tem suavidade $m$ se $\int (D^m r(x))^2\,dx \le K_1$
        ($m$ derivadas de quadrado integrável). Quanto maior $m$, mais suave.
\end{description}

%==============================================================================
\section{A taxa de convergência do KNN}
%==============================================================================
Considere o estimador KNN $\rhat_k(\x)=\frac1k\sum_{i\in\mathcal{N}_{\x}}Y_i$
(Aula~04). Sob suavidade Lipschitz, conseguimos limitar seu risco.

\begin{lema}[Distância média aos $k$ vizinhos]\label{lem:viz}
Se $\X$ é limitado e $p\ge 3$, então
\[
  \E\!\left[\Big(\frac1k\sum_{i\in\mathcal{N}_{\X}}\norm{\X_i-\X}\Big)^2\right]
  \le C\,\Big(\frac{k}{n}\Big)^{2/p}.
\]
\end{lema}
Em palavras: quanto mais vizinhos ($k$) eu exijo, mais \emph{longe} eles ficam --- e
essa distância cresce com $(k/n)^{1/p}$. O expoente $1/p$ é o gérmen da maldição da
dimensionalidade.

\begin{teorema}[Risco do KNN]\label{teo:knn}
Sob as hipóteses do Lema~\ref{lem:viz}, com $r$ $L$-Lipschitz e
$\sup_{\x}\Var(Y\mid\x)\le\sigma^2<\infty$,
\[
  \E\big[(\rhat_k(\X)-r(\X))^2\big]
   \;\le\; K\,\Big(\frac{k}{n}\Big)^{2/p} + \frac{\sigma^2}{k}.
\]
Otimizando em $k$ (tomando $k \asymp n^{2/(2+p)}$), obtém-se
\[
  \E\big[(\rhat_k(\X)-r(\X))^2\big] \;\le\; K'\, n^{-2/(2+p)}.
\]
\end{teorema}
\begin{proof}
Condicione nas covariáveis $\widetilde{\x}=(\x_1,\dots,\x_n,\x)$. Pela decomposição
viés--variância (Aula~01),
\[
  \E\big[(\rhat_k(\x)-r(\x))^2\mid\widetilde{\x}\big]
  = \underbrace{\Big(\tfrac1k\!\sum_{i\in\mathcal{N}_{\x}}\! r(\x_i)-r(\x)\Big)^2}_{\text{viés}^2}
  + \underbrace{\Var(\rhat_k(\x)\mid\widetilde{\x})}_{\text{variância}}.
\]
\emph{Viés:} como $r$ é $L$-Lipschitz,
$\big|\tfrac1k\sum_{i\in\mathcal{N}_{\x}}r(\x_i)-r(\x)\big|
\le \tfrac{L}{k}\sum_{i\in\mathcal{N}_{\x}}\norm{\x_i-\x}$.
\emph{Variância:} as respostas são independentes dadas as covariáveis e cada uma
tem variância $\le\sigma^2$; a média de $k$ delas tem variância
$\Var(\rhat_k(\x)\mid\widetilde{\x})\le \sigma^2/k$. Logo
\[
  \E\big[(\rhat_k(\x)-r(\x))^2\mid\widetilde{\x}\big]
  \le \Big(\tfrac{L}{k}\!\sum_{i\in\mathcal{N}_{\x}}\!\norm{\x_i-\x}\Big)^2 + \frac{\sigma^2}{k}.
\]
Tomando esperança e aplicando o Lema~\ref{lem:viz} ao primeiro termo, resulta
$\E[(\rhat_k(\X)-r(\X))^2]\le K(k/n)^{2/p}+\sigma^2/k$. Para minimizar em $k$,
iguale (a menos de constantes) os dois termos: $(k/n)^{2/p}\asymp 1/k$, isto é,
$k^{1+2/p}\asymp n^{2/p}$, ou seja $k\asymp n^{2/(2+p)}$, e substituindo obtém-se
a taxa $n^{-2/(2+p)}$.
\end{proof}

\begin{ideia}
A prova exibe explicitamente o balanço viés--variância da Aula~01: o termo
$(k/n)^{2/p}$ \emph{cresce} com $k$ (viés) e $\sigma^2/k$ \emph{decresce} com $k$
(variância). O $k$ ótimo equilibra os dois --- e a validação cruzada (Aula~03) é a
maneira prática de encontrá-lo.
\end{ideia}

\subsection{Séries ortogonais e taxa minimax}
Para séries ortogonais (base de Fourier) com suavidade de Sobolev $m$, um argumento
análogo (viés do truncamento em $J$ termos $\le K_1 J^{-2m}$; variância $\le K_2 J/n$)
dá
\[
  \E\big[(\rhat_J(\X)-r(\X))^2\big] \le K_1 J^{-2m} + K_2\frac{J}{n},
  \quad\text{e com } J\asymp n^{1/(2m+1)}:\quad
  \E[\cdots] \le K\, n^{-2m/(2m+1)}.
\]
Note que mais suavidade ($m$ maior) $\Rightarrow$ taxa mais rápida.

\begin{observacao}[Taxa minimax]
Essas taxas não são apenas \emph{limitantes superiores} de um método específico:
para a classe $\mathcal{S}$ de funções suaves elas coincidem com a \textbf{taxa
minimax}
\[
  \inf_{\rhat}\ \sup_{r\in\mathcal{S}}\ \E\big[(\rhat(\X)-r(\X))^2\big]
  \asymp n^{-\alpha/(\alpha+p)},
\]
ou seja, \emph{nenhum} estimador faz essencialmente melhor no pior caso. Um
estimador que atinge essa taxa é dito \emph{minimax-ótimo}. ($\alpha$ liga-se à
suavidade.)
\end{observacao}

%==============================================================================
\section{A maldição da dimensionalidade}
%==============================================================================
Olhe a taxa minimax $n^{-\alpha/(\alpha+p)}$: para $p$ grande, o expoente
$\alpha/(\alpha+p)$ se aproxima de $0$ --- a convergência fica \emph{lentíssima}.

\begin{teorema}[Tamanho de amostra explode com $p$]
Para garantir risco $\le\delta$ na classe $\mathcal{S}$, é necessário
\[
  n \;\gtrsim\; \Big(\frac{K}{\delta}\Big)^{(\alpha+p)/\alpha},
\]
um crescimento \textbf{exponencial} em $p$.
\end{teorema}
Mesmo dimensões moderadas exigem amostras astronômicas. Esta é a
\textbf{maldição da dimensionalidade} (Bellman, 1966).

\begin{ideia}[A intuição da vizinhança vazia]
Para estimar $r(\x)$ por média local, preciso de pontos \emph{perto} de $\x$. Em
$\R^p$ com $p$ grande, qualquer vizinhança de raio fixo tem volume
desprezível: com alta probabilidade \emph{não há nenhum dado dentro dela}. Para
capturar $k$ vizinhos preciso de um raio $\asymp (k/n)^{1/p}$, que cresce com $p$
--- os ``vizinhos'' deixam de ser próximos, e o viés explode (é exatamente o
Lema~\ref{lem:viz}).
\end{ideia}

\begin{atencao}
A maldição diz que, \textbf{sem suposições extras}, estimação não paramétrica é
impossível em alta dimensão para $n$ realista. A saída não é um método melhor ---
é uma \emph{estrutura} adicional no problema. As duas mais úteis: esparsidade e
redundância.
\end{atencao}

%==============================================================================
\section{Escapando da maldição}
%==============================================================================
\subsection{Esparsidade}
A hipótese de \textbf{esparsidade} diz que $r$ depende apenas de um subconjunto
$S\subset\{1,\dots,p\}$ de covariáveis:
\[
  r(\x) = r\big((x_i)_{i\in S}\big),\qquad \abs{S}=s \ll p.
\]
Embora haja $p$ covariáveis, só $s$ importam. Sob esparsidade a dimensão
\emph{efetiva} é $s$, e a taxa melhora para $\approx n^{-\alpha/(\alpha+s)}$.

\begin{emsala}
Eis a resposta da pergunta de abertura (Amazon): no problema de texto, a maioria
das palavras é irrelevante para a nota --- problema \emph{esparso}. Métodos que
\textbf{selecionam variáveis} (Lasso, árvores/florestas, XGBoost) prosperam; o KNN,
que usa \emph{todas} as covariáveis na distância com peso igual, é arrastado pelas
irrelevantes e vai mal. KNN \emph{não} faz seleção de variáveis.
\end{emsala}

\begin{observacao}
Métodos que se adaptam à esparsidade: \emph{Lasso} (Aula~02), \emph{florestas
aleatórias} (cuja taxa, por Biau 2012, depende só do número de variáveis
relevantes, pois divisões inúteis raramente são escolhidas) e \emph{SpAM}
(modelos aditivos esparsos, que combinam aditividade com $\ell_1$).
\end{observacao}

\subsection{Redundância}
A hipótese de \textbf{redundância} diz que as covariáveis, embora em número $p$,
vivem (aproximadamente) sobre uma \emph{subvariedade} de dimensão intrínseca
$u\ll p$. Exemplos:
\begin{itemize}[leftmargin=1.4cm,nosep]
  \item altura, peso e IMC: três variáveis, mas duas determinam a terceira;
  \item imagens: \emph{pixels} adjacentes são quase iguais --- enorme redundância.
\end{itemize}
Sob redundância, KNN e regressão local \emph{se adaptam automaticamente} à dimensão
intrínseca: sua taxa passa a ser $n^{-4/(4+u)}$, muito mais rápida que
$n^{-4/(4+p)}$ quando $u\ll p$. SVR com kernel gaussiano e séries \emph{espectrais}
(base construída a partir dos dados) também exploram redundância.

\begin{ideia}
Esparsidade $\neq$ redundância. \emph{Esparsidade}: poucas covariáveis importam
(seleção de variáveis ajuda). \emph{Redundância}: muitas covariáveis, mas
correlacionadas/numa subvariedade (redução de dimensionalidade ajuda --- tema da
Aula extra de PCA). Métodos diferentes exploram cada uma.
\end{ideia}

%==============================================================================
\section{Síntese prática}
%==============================================================================
\begin{center}\small
\renewcommand{\arraystretch}{1.3}
\begin{tabular}{@{}lll@{}}
\toprule
\textbf{Situação} & \textbf{O que ajuda} & \textbf{Métodos} \\
\midrule
$p$ pequeno, $r$ suave & nada extra & KNN, NW, regressão local \\
Poucas covariáveis relevantes & seleção de variáveis & Lasso, florestas, XGBoost, SpAM \\
Covariáveis redundantes & dimensão intrínseca & KNN/local, SVR, séries espectrais, PCA \\
\bottomrule
\end{tabular}
\end{center}

O notebook \texttt{Aula prática 05} mede, por simulação, cada passo do argumento
desta aula: o expoente $1/p$ do Lema, a taxa do Teorema, e a diferença entre
esparsidade e redundância.

\begin{atencao}
Mensagem para levar para casa: \textbf{não existe almoço grátis}. A escolha do
método deve refletir a \emph{estrutura} que você acredita existir nos dados. A
teoria desta aula é o que torna essa escolha informada em vez de adivinhação.
\end{atencao}

%==============================================================================
\section*{Resumo}
%==============================================================================
\begin{itemize}[leftmargin=1.4cm]
  \item Taxas de convergência medem a qualidade assintótica; só existem sob
        \emph{suavidade} (Lipschitz, Sobolev).
  \item KNN: risco $\le K(k/n)^{2/p}+\sigma^2/k$, ótimo em
        $k\asymp n^{2/(2+p)}$, taxa $n^{-2/(2+p)}$ (Teo.~\ref{teo:knn}).
  \item Séries ($m$-suaves): taxa $n^{-2m/(2m+1)}$; ambas atingem a taxa
        \emph{minimax} $n^{-\alpha/(\alpha+p)}$.
  \item \textbf{Maldição da dimensionalidade}: a taxa degrada com $p$ e o $n$
        necessário cresce exponencialmente.
  \item \emph{Esparsidade} (poucas relevantes) $\to$ Lasso/florestas;
        \emph{redundância} (dimensão intrínseca baixa) $\to$ KNN/local, séries
        espectrais, PCA.
\end{itemize}

\paragraph{Leitura recomendada.} [AME] §4.13 (taxas do KNN e séries ortogonais),
Capítulo~5 (§5.1 maldição da dimensionalidade, §5.1.1 esparsidade, §5.1.2
redundância, §5.5--5.6 florestas e SpAM).

\end{document}
